\subsection{محدودیت‌ها و تحلیل شکست}
محدودیت‌ها و عوامل موثر بر افت دقت در آزمون‌ها عبارتند از:
\begin{itemize}
    \item \textbf{افت عملکرد به سطح \lr{PASS\_LOW} در مدل‌های زبانی عمیق (\lr{RoBERTa} و \lr{GPT-2}):} اعمال یکنواخت تنک‌سازی به میزان $50\%$ بدون تنظیم مجدد آداپتورهای \lr{LoRA} در طول آزمون استنتاج صفر-شات، موجب حذف بخش‌هایی از اتصالات حیاتی در لایه‌های میانی شبکه شده که منجر به تنزل دقت در دسته‌بندی توالی‌های ظریف شده است.
    \item \textbf{نیاز به بازیابی آموزش (\lr{Fine-Tuning Recovery}):} بر خلاف روش‌های فاکتورگیری که هندسه فضای پیوسته را حفظ می‌کنند، هرس ساختاریافته ماتریس‌ها را مقید به صفرسازی کامل اتصالات می‌کند؛ لذا برای احیای دقت کامل مدل‌هایی نظیر \lr{GPT-2}، آموزش تکمیلی گام‌به‌گام ضروری خواهد بود.
\end{itemize}
